\unchapter{Conclusions}

The central aim of this thesis was to evaluate how well information
systems are integrated within Project Management Office environments
in medium and large Latvian enterprises, and to explore whether this
integration maturity has any meaningful connection to reporting
automation, data quality, and governance standardization. This
concluding section confirms the completion of all research tasks,
reflects on the four hypotheses, and assesses whether the overall
aim of the thesis has been achieved.

\textbf{Confirmation of research tasks.}

Task 1 was addressed in Chapter~1, where a thorough review of
existing literature was carried out covering PMO digitalization,
information system integration architectures, and maturity models
applicable to IS environments. This review laid the conceptual and
theoretical groundwork for the entire study.

Task 2 was also completed in Chapter~1, where a four-level
Integration Maturity Model was constructed and made measurable
through clearly defined criteria at each level, encompassing
integration mechanisms, degree of automation, reporting structure,
and the extent of BI centralization.

Task 3 was fulfilled in Chapter~2, which documented the development
and validation of a structured survey instrument. The instrument
included both Likert-scale and categorical items, each directly
tied to the dimensions of the Integration Maturity Model.

Task 4 involved distributing the survey to representatives of
medium and large Latvian enterprises employing 50 or more people
and maintaining a formal PMO structure. In total, 40 usable
responses were received, which fell within the 30--50 target
range established in the methodology.

Task 5 was carried out in Chapter~3, where a structured expert
interview was held with a seasoned project management professional.
The interview served to enrich the quantitative survey findings
with qualitative perspectives and practical architectural insights.

Task 6 was also completed in Chapter~3, where descriptive
statistical analysis was used to map out how integration maturity
levels were distributed across the organizations in the sample.

Task 7 was fulfilled in Chapter~3 through the application of
Spearman correlation and simple regression analysis, which were
used to investigate the relationships between integration maturity
and the three dependent variables: reporting automation, data
accuracy, and governance standardization.

Task 8 was addressed in Chapter~4, where the empirical results
were interpreted, conclusions were drawn in relation to all four
hypotheses, and five concrete recommendations were put forward
to help Latvian enterprises strengthen their PMO IS integration
maturity.

\textbf{Validation of hypotheses.}

All four hypotheses received support from the empirical data
gathered in this study.

Hypothesis 1 was confirmed. The survey data showed that 70\%
of the sampled PMO environments are operating at intermediate
integration maturity levels — specifically Level 2 (Developing)
or Level 3 (Defined) — rather than at the most advanced stage.
With only 15\% of organizations having reached Level 4, it is
clear that the majority of Latvian PMOs remain in a transitional
phase of IS integration development.

Hypothesis 2 was confirmed. A statistically significant positive
correlation was identified between integration maturity and
reporting automation (r=0.405, p=0.010). This finding indicates
that organizations operating at higher integration maturity levels
consistently achieve greater degrees of automated reporting,
thereby reducing their dependence on manual data consolidation.

Hypothesis 3 was confirmed. A statistically significant positive
correlation was found between integration maturity and data
quality (r=0.321, p=0.043). Although this was the weakest
association among the three, it nonetheless confirms that
organizations with better integrated systems tend to perceive
their data as more accurate and accessible in a timely manner.

Hypothesis 4 was confirmed. A statistically significant positive
correlation was established between integration maturity and
governance standardization (r=0.403, p=0.010). Organizations
with higher integration maturity exhibited noticeably stronger
governance standardization, a finding further reinforced by the
expert interview, which highlighted the transformative role of
centralized BI dashboards in bringing consistency to governance
processes.

\textbf{Achievement of the aim.}

The aim of this thesis has been fully achieved. Through a
combination of quantitative survey analysis and qualitative
expert interview findings, the study successfully assessed IS
integration maturity across PMO environments in medium and large
Latvian enterprises. The results clearly demonstrate statistically
significant relationships between integration maturity and all
three dependent variables. It can be concluded that most Latvian
PMO environments currently operate at intermediate maturity levels,
and that organizations with higher integration maturity
consistently perform better in terms of reporting automation,
data quality, and governance standardization. The aim of the
thesis has therefore been reached.

\textbf{Practical and theoretical value.}

This study offers value on both a practical and theoretical level.
From a practical standpoint, five evidence-based recommendations
were developed to guide Latvian enterprises in improving their
PMO IS integration maturity. These cover a range of actionable
steps: building API-based integration between core systems,
adopting centralized BI dashboards, exploring low-code platforms
as a cost-effective alternative for medium-sized organizations,
resolving compatibility issues between local Latvian accounting
software and international PM tools, and investing in data
literacy programmes for project management staff.

From a theoretical standpoint, this study delivers the first
empirical validation of a PMO-specific IS Integration Maturity
Model within the Latvian enterprise context. It demonstrates
that maturity modeling can be meaningfully applied to IS
integration in PMO settings, and adds a new layer of empirical
evidence to the growing body of scholarship on digital
transformation in Baltic and Eastern European organizations.

All eight tasks outlined in the Introduction were successfully
completed. Tasks 1 and 2 were addressed in Chapter~1, Task 3
in Chapter~2, Tasks 4--7 in Chapter~3, and Task 8 in Chapter~4.
All four research hypotheses were validated by the empirical
data, and the aim of the thesis has been fully achieved.
